\section{1. Introduction}

Light field imaging captures both spatial and angular information of light rays, offering significant advantages in 3D reconstruction, autonomous navigation, and robotic vision [1]. Depth estimation serves as a core task in this domain, extracting geometric structures from multi-view consistency cues [2]. Commercial light field cameras, such as Lytro and Raytrix, enable dense angular sampling, facilitating accurate disparity calculation through epipolar plane image (EPI)s (EPIs) [3]. Consequently, numerous learning-based and traditional frameworks have achieved notable accuracy in controlled environments [4].

However, obtaining the accurate shape of glossy, transparent, or metallic objects remains a hard challenge for both classical algorithms and deep networks [5]. Most current depth estimation methods rely heavily on the Lambertian assumption, where surface radiance remains constant across different viewpoints [6]. In real-world scenarios, surfaces frequently exhibit non-Lambertian reflectance, including specular reflections and transparent refractions. These complex optical phenomena violate the photo-consistency principle, causing severe distortions and discontinuities in EPI slope patterns [7]. As a result, standard matching costs and learning-based feature extractors fail to establish reliable correspondences across angular views.

The performance degradation in non-Lambertian and mixed-material scenes is substantial. While state-of-the-art models achieve a mean absolute error (MAE) below 0.05 on purely Lambertian datasets, their errors often exceed 0.30 on urban mixed scenes [8]. This represents a severe performance degradation of over 500\% when encountering complex reflectance. Furthermore, real-world light field datasets exhibit a severe long-tail distribution regarding material types . Minority domains, such as pure non-Lambertian scenes, possess limited training samples compared to dominant Lambertian backgrounds. During joint training, the gradients from majority domains easily overwhelm those from minority domains. This prevents the network from learning robust material-invariant features and severely restricts generalization.

Addressing these challenges requires a unified physical model that explicitly handles diverse reflectance properties without relying on isolated heuristics. Moreover, algorithm design must respect the physical limitations of current light field hardware. For instance, mainstream cameras provide a $9 \times 9$ angular resolution, which strictly limits the Nyquist frequency in the angular domain. Attempting to classify materials using angular frequency analysis under such sparse sampling is physically unfeasible. Our empirical analysis reveals that the non-DC spectral energy remains below 3\% under this configuration. Therefore, developing a robust framework necessitates advanced architectural designs and a rigorous understanding of hardware-induced physical boundaries. Establishing such a unified baseline is essential for advancing light field perception in unconstrained environments.

...increase significantly when applied to non-Lambertian and mixed-material scenes. To understand this degradation, we examine the limitations of existing approaches across different developmental stages.

Traditional light field depth estimation methods primarily rely on handcrafted matching costs and epipolar plane image (EPI) (EPI) slope analysis [4]. These approaches explicitly assume Lambertian reflectance to enforce photo-consistency across angular views [6]. However, this assumption fails on glossy or transparent surfaces, where specular highlights and refractions cause severe EPI distortions [9]. Consequently, traditional algorithms produce fragmented depth maps in non-Lambertian regions. Furthermore, they lack adaptive mechanisms for mixed-material scenes, requiring tedious manual priors to separate diffuse and specular components [10].

To overcome the limitations of handcrafted features, early deep learning methods employ convolutional neural networks to extract spatial-angular representations [11]. While these models improve accuracy on synthetic Lambertian datasets, they implicitly learn feature mappings without physical interpretability [12]. As a result, they struggle to explicitly decouple diffuse and specular reflections in complex environments. Moreover, these networks are typically trained on single-domain datasets with uniform sampling strategies. When applied to multi-domain scenarios with long-tailed distributions, the gradients from dominant Lambertian scenes overwhelm minority non-Lambertian domains [13]. This gradient domination severely degrades the generalization performance on data-scarce reflective surfaces.

Recent advanced architectures integrate attention mechanisms, multi-scale processing, and frequency-domain analysis to tackle complex reflectance [14]. Some studies hypothesize that different materials exhibit distinct spectral distributions, proposing angular 2D fast Fourier transform (FFT) for material classification [15]. However, this approach overlooks the physical sampling limits of commercial light field cameras. Under standard $9 \times 9$ angular resolutions, the Nyquist frequency restricts the effective frequency bins, rendering angular frequency analysis practically infeasible. Additionally, these state-of-the-art models often incur substantial parameter overhead and lack a unified baseline for diverse scenes [16]. They also fail to identify the theoretical performance ceiling of pure EPI-based architectures in highly textured environments.

Therefore, there is an urgent need for a unified, physically grounded, and lightweight framework that explicitly models non-Lambertian phenomena. Such a method must overcome domain imbalance in multi-scene training and respect hardware sampling limits for robust depth estimation.

In this paper, we propose a unified dual-mask physical model for light field depth estimation that explicitly decouples diffuse and specular reflections. To address the failure of the Lambertian assumption in complex scenes, our framework integrates a dual-mask modeling mechanism with a multi-directional epipolar plane image (EPI) (EPI) processing network. Specifically, the model employs an angular fast Fourier transform module to generate medium and angular direction masks, which subsequently guide the adaptive fusion of pure defocus and angular frequency depth features. This compact architecture, containing only 754K parameters, is trained end-to-end at a full resolution of $384 \times 384$ to simultaneously handle Lambertian, non-Lambertian, and mixed urban scenarios.

Our main contributions are summarized as follows:

\begin{itemize}
  \item \textbf{Contribution 1}: We develop a unified dual-mask physical architecture that integrates medium and angular direction masks to decouple material-specific depth cues. By combining a pure defocus branch for non-Lambertian regions and an angular frequency branch for Lambertian areas, our model resolves the inability of traditional methods to process mixed materials. This lightweight design achieves a robust baseline, yielding an overall mean absolute error (MAE) of 0.1795 across diverse complex scenes.
  \item \textbf{Contribution 2}: We establish the physical limits of angular frequency domain material classification by rigorously falsifying its feasibility under current $9 \times 9$ angular resolutions. Through four independent spectral analysis methods and Nyquist frequency evaluations, we demonstrate that non-DC energy is less than 3\%, proving that at least a $17 \times 17$ angular sampling is required. This negative finding provides explicit theoretical guidance for future light field sensor design and prevents blind exploration in low-resolution frequency analysis.
  \item \textbf{Contribution 3}: We introduce a domain-balanced sampling strategy and resolution scaling to mitigate gradient domination in multi-domain long-tailed distributions. By oversampling minority domains (e.g., weight 39.0 for non-Lambertian) and increasing the input resolution to $384 \times 384$, we substantially enhance generalization in data-scarce regions. Furthermore, we diagnose the architectural limits of EPI-based networks in complex textures, indicating that surpassing current performance ceilings necessitates non-EPI or hybrid paradigms.
\end{itemize}
Extensive experiments on four benchmark datasets demonstrate the effectiveness of our proposed method. The unified model achieves an overall MAE of 0.1795, with specific errors of 0.152 for Lambertian, 0.211 for mixed urban, and 0.235 for non-Lambertian scenes. These results validate that integrating physical masks with multi-directional EPI features significantly improves depth estimation accuracy in challenging non-Lambertian environments.